\documentclass[a4paper]{article}

\usepackage{graphicx}
\usepackage{amsmath,amssymb}
\usepackage{minted}
\usepackage{multirow}
\usepackage{float}
\usepackage{algorithm}
\usepackage{algpseudocode}
\usepackage{todonotes}
\usepackage{forest}
\usepackage{xstring}
\usepackage{xcolor}
\usepackage[most]{tcolorbox}
\usepackage{xparse}
\usepackage{natbib}
\usepackage{adjustbox}
\usepackage[active,tightpage]{preview}

\PreviewBorder=4pt

\usetikzlibrary{graphs,graphdrawing,arrows.meta,positioning,calc,fit,shapes.geometric}
\usegdlibrary{layered}

\input{/nd_language_ai_project/src/nd_language_ai/formal/computer/programming/tex/latex/code_clips/external_graphviz.tex}
\input{/nd_language_ai_project/src/nd_language_ai/formal/computer/programming/tex/latex/part/control_sequence/kind/semtag/semtag.tex}
\input{/nd_language_ai_project/src/nd_language_ai/formal/computer/programming/tex/latex/code_clips/seehere.tex}
\input{/nd_language_ai_project/src/nd_language_ai/formal/computer/programming/tex/latex/code_clips/deep_structure.tex}

\title{TikZ Pencil Connectors, Connection Points, and Layout Directions}
\date{2026-09-21}
\author{Mohammad Rahmani}

\begin{document}
   \maketitle

   \begin{preview}
      \section{Connection-point workflow}
      Rectangle and ellipse nodes expose four perimeter connection points by
      default.  Additional points are added from the node long-press menu and
      are distributed at equal perimeter distance.  Their counts are persisted
      in compile-safe TikZ comments.

      \section{Apple Pencil connector workflow}
      A Pencil stroke starting at a connection point can begin while Select is
      active.  Mind Mirror reactivates the most recently chosen connector kind
      and previews that connector.  Releasing over a connection point on
      another node creates the edge.  Releasing on blank canvas opens a choice
      to create and connect to a rectangle, create and connect to an ellipse,
      or discard the connector.

      \section{Automatic layout}
      The TikZ options menu now exposes Top to Bottom, Bottom to Top, Left to
      Right, Right to Left, Right Tree, and Left Tree layouts.
   \end{preview}

   \bibliographystyle{plainnat}
   \bibliography{/home/donkarlo/Dropbox/repo/nd_shared_research/bibliography/bibliography.bib}
\end{document}
